

\section{周二}
\subsection{希利尼人見耶穌 (約翰 12:20-22)}

那時，上來過節禮拜的人中，有幾個希利尼人。他們來見加利利伯賽大的腓力，求他說：先生，我們願意見耶穌。腓力去告訴安得烈，安得烈同腓力去告訴耶穌。
\subsection{最後的公開講論（約12:23-50）}

\begin{table}[htbp]
  \centering
    \begin{tabular}{|p{3em}|p{38.645em}|}
    \toprule
    言者    & 講話內容 \\
    \midrule
    耶穌    & 人子得榮耀的時候到了。….父啊，願你榮耀你的名！ \\
    \midrule
    父神    & 當時就有聲音從天上來說：「我已經榮耀了我的名，還要再榮耀。」 \\
    \midrule
    眾人    & 站在旁邊的眾人聽見，就說：「打雷了。」 \\
    \midrule
    有人    & 還有人說：「有天使對他說話。」 \\
    \midrule
    耶穌    & 這聲音不是為我，是為你們來的。現在這世界受審判，這世界的王要被趕出去。我若從地上被舉起來，就要吸引萬人來歸我。 \\
    \midrule
    眾人    & 我們聽見律法上有話說基督是永存的，你怎麼說人子必須被舉起來呢？這人子是誰呢？ \\
    \midrule
    耶穌    & 光在你們中間還有不多的時候，應當趁著有光行走，免得黑暗臨到你們。那在黑暗裡行走的，不知道往何處去。你們應當趁著有光，信從這光，使你們成為光明之子。 \\
    \midrule
    官長    & 官長中卻有好些信他的，只因法利賽人的緣故，就不承認，恐怕被趕出會堂。這是因他們愛人的榮耀過於愛神的榮耀。 \\
    \midrule
    耶穌    & 信我的，不是信我，乃是信那差我來的… \\
    \bottomrule
    \end{tabular}%
    \caption{最後的公開講論}
  \label{tab:addlabel}%
\end{table}


耶穌說：人子得榮耀的時候到了。我實實在在的告訴你們，一粒麥子不落在地裡死了，仍舊是一粒，若是死了，就結出許多子粒來。愛惜自己生命的，就失喪生命；在這世上恨惡自己生命的，就要保守生命到永生。若有人服事我，就當跟從我；我在哪裡，服事我的人也要在那裡；若有人服事我，我父必尊重他。我現在心裡憂愁，我說什麼才好呢？父啊，救我脫離這時候；但我原是為這時候來的。父啊，願你榮耀你的名！當時就有聲音從天上來，說：我已經榮耀了我的名，還要再榮耀。站在旁邊的眾人聽見，就說：打雷了。還有人說：有天使對他說話。耶穌說：這聲音不是為我，是為你們來的。現在這世界受審判，這世界的王要被趕出去。我若從地上被舉起來，就要吸引萬人來歸我。耶穌這話原是指著自己將要怎樣死說的。眾人回答說：我們聽見律法上有話說，基督是永存的，你怎麼說人子必須被舉起來呢？這人子是誰呢？耶穌對他們說：光在你們中間還有不多的時候，應當趁著有光行走，免得黑暗臨到你們；那在黑暗裡行走的，不知道往何處去。你們應當趁著有光，信從這光，使你們成為光明之子。耶穌說了這話，就離開他們隱藏了。他雖然在他們面前行了許多神蹟，他們還是不信他。這是要應驗先知以賽亞的話，說：主啊，我們所傳的有誰信呢？主的膀臂向誰顯露呢？他們所以不能信，因為以賽亞又說：主叫他們瞎了眼，硬了心，免得他們眼睛看見，心裡明白，回轉過來，我就醫治他們。以賽亞因為看見他的榮耀，就指著他說這話。雖然如此，官長中卻有好些信他的，只因法利賽人的緣故，就不承認，恐怕被趕出會堂。這是因他們愛人的榮耀過於愛神的榮耀。耶穌大聲說：信我的，不是信我，乃是信那差我來的。人看見我，就是看見那差我來的。我到世上來，乃是光，叫凡信我的，不住在黑暗裡。若有人聽見我的話不遵守，我不審判他。我來本不是要審判世界，乃是要拯救世界。棄絕我、不領受我話的人，有審判他的─就是我所講的道在末日要審判他。因為我沒有憑著自己講，惟有差我來的父已經給我命令，叫我說什麼，講什麼。我也知道他的命令就是永生。故此，我所講的話正是照著父對我所說的。
\subsection{在橄欖山論將來之事} 
\subsubsection{末了的預兆-太24:1-14/可13：1-13/路21：5-19}
\paragraph{馬太 24:1-14}
耶穌出了聖殿，正走的時候，門徒進前來，把殿宇指給他看。
耶穌對他們說：你們不是看見這殿宇嗎？我實在告訴你們，將來在這裡沒有一塊石頭留在石頭上，不被拆毀了。
耶穌在橄欖山上坐著，門徒暗暗的來說：請告訴我們，什麼時候有這些事？你降臨和世界的末了有什麼預兆呢？
耶穌回答說：你們要謹慎，免得有人迷惑你們。
因為將來有好些人冒我的名來，說：我是基督，並且要迷惑許多人。
你們也要聽見打仗和打仗的風聲，總不要驚慌；因為這些事是必須有的，只是末期還沒有到。
民要攻打民，國要攻打國；多處必有饑荒、地震。
這都是災難（災難：原文是生產之難）的起頭。
那時，人要把你們陷在患難裡，也要殺害你們；你們又要為我的名被萬民恨惡。
那時，必有許多人跌倒，也要彼此陷害，彼此恨惡；
且有好些假先知起來，迷惑多人。
只因不法的事增多，許多人的愛心才漸漸冷淡了。
惟有忍耐到底的，必然得救。
這天國的福音要傳遍天下，對萬民作見證，然後末期才來到。

\paragraph{馬可 13：1-13}
耶穌從殿裡出來的時候，有一個門徒對他說：夫子，請看，這是何等的石頭！何等的殿宇！
耶穌對他說：你看見這大殿宇嗎？將來在這裡沒有一塊石頭留在石頭上，不被拆毀了。
耶穌在橄欖山上對聖殿而坐，彼得、雅各、約翰，和安得烈暗暗的問他說：
請告訴我們，什麼時候有這些事呢？這一切事將成的時候有什麼預兆呢？
耶穌說：你們要謹慎，免得有人迷惑你們。
將來有好些人冒我的名來，說：我是基督，並且要迷惑許多人。
你們聽見打仗和打仗的風聲，不要驚慌。這些事是必須有的，只是末期還沒有到。
民要攻打民，國要攻打國；多處必有地震、饑荒。這都是災難（災難：原文是生產之難）的起頭。
但你們要謹慎；因為人要把你們交給公會，並且你們在會堂裡要受鞭打，又為我的緣故站在諸侯與君王面前，對他們作見證。
然而，福音必須先傳給萬民。
人把你們拉去交官的時候，不要預先思慮說什麼；到那時候，賜給你們什麼話，你們就說什麼；因為說話的不是你們，乃是聖靈。
弟兄要把弟兄，父親要把兒子，送到死地；兒女要起來與父母為敵，害死他們；
並且你們要為我的名被眾人恨惡。惟有忍耐到底的，必然得救。


\paragraph{路加 21：5-19}
有人談論聖殿是用美石和供物妝飾的；
耶穌就說：論到你們所看見的這一切，將來日子到了，在這裡沒有一塊石頭留在石頭上，不被拆毀了。
他們問他說：夫子！什麼時候有這事呢？這事將到的時候有什麼預兆呢？
耶穌說：你們要謹慎，不要受迷惑；因為將來有好些人冒我的名來，說：我是基督，又說：時候近了，你們不要跟從他們！
你們聽見打仗和擾亂的事，不要驚惶；因為這些事必須先有，只是末期不能立時就到。
當時，耶穌對他們說：民要攻打民，國要攻打國；
地要大大震動，多處必有饑荒、瘟疫，又有可怕的異象和大神蹟從天上顯現。
但這一切的事以先，人要下手拿住你們，逼迫你們，把你們交給會堂，並且收在監裡，又為我的名拉你們到君王諸侯面前。
但這些事終必為你們的見證。
所以，你們當立定心意，不要預先思想怎樣分訴；
因為我必賜你們口才、智慧，是你們一切敵人所敵不住、駁不倒的。
連你們的父母、弟兄、親族、朋友也要把你們交官；你們也有被他們害死的。
你們要為我的名被眾人恨惡，
然而，你們連一根頭髮也必不損壞。
你們常存忍耐，就必保全靈魂（或作：必得生命）。
\paragraph{综合}
耶穌從殿裡出來的時候，正走的時候，有些人谈论到圣殿，说它是用精美的石头和供奉的礼物装饰的。有一個門徒把殿宇指給他看對他說：「夫子，請看，這是何等的石頭，何等的殿宇！」耶穌對他說：「你看見這大殿宇嗎？我實在告訴你們，將來在這裡沒有一塊石頭留在石頭上不被拆毀了。」
耶穌在橄欖山上對聖殿而坐，彼得、雅各、約翰和安得烈暗暗地問他說： 4 「請告訴我們，什麼時候有這些事呢？這一切事將成的時候，你降臨和世界的末了有什麼預兆呢？」 5 耶穌說：「你們要謹慎，免得有人迷惑你們。 6 將來有好些人冒我的名來，說『我是基督』，並且要迷惑許多人。又说‘时候快到了。’你们不要跟从他们。 7 你們也要聽見打仗和打仗的風聲，要注意，不要驚慌，這些事是必須有的，只是末期(結局)還沒有到。8 民要攻打民，國要攻打國，一個民族要起來反對另一個民族，一個國家要起來反對另一個國家. 多處(到處)必有地震、饑荒和瘟疫，這都是災難 a的起頭。這些都像臨產陣痛的開始。9 会有恐怖的景象和天上来的大征兆。 12 不过在所有这些事以前，他们将对你们下手，你們要謹慎，因為那時，他們將要出賣你們，人要把你們陷在患難裡，也要殺害你們；你們又要為我的名被萬民恨惡。人要把你們交給公會(議會)，送进监狱. 並且你們在會堂裡要受鞭打，又為我的緣故站在諸侯與君王面前，對他們作見證。 10 然而，福音必須先傳給萬民。 11 人把你們拉去交官的時候，应当心里镇定，不要預先思慮說什麼，到那時候，賜給你們什麼話，你們就說什麼；因為說話的不是你們，乃是聖靈。是所有反对你们的人都不能抵挡、不能驳倒的。 你们甚至将被父母、兄弟、亲戚和朋友出卖，他们还会害死你们当中的一些人。12 弟兄要把弟兄，父親要把兒子，送到死地；兒女要起來與父母為敵，害死他們。那時，許多人將被絆倒，彼此出賣，彼此憎恨； 11 並且會有許多假先知起來，迷惑許多人； 12 由於罪惡增多，許多人的愛心就會冷淡。 13 但是忍耐到底的，這個人將會得救。 14 這天國的福音將被傳遍天下，好對萬國做見證，然後結局才會到來。 13 並且你們要為我的名被眾人恨惡，可是你们连一根头发也绝不会失去。唯有忍耐到底的，必然得救。


\subsubsection{耶路撒冷的大災難-太24:15-28/可13：14-23/路21：20-24}
\paragraph{馬太 24:15-28}
你們看見先知但以理所說的那行毀壞可憎的站在聖地（讀這經的人須要會意）。
那時，在猶太的，應當逃到山上；
在房上的，不要下來拿家裡的東西；
在田裡的，也不要回去取衣裳。
當那些日子，懷孕的和奶孩子的有禍了。
你們應當祈求，叫你們逃走的時候，不遇見冬天或是安息日。
因為那時必有大災難，從世界的起頭直到如今，沒有這樣的災難，後來也必沒有。
若不減少那日子，凡有血氣的總沒有一個得救的；只是為選民，那日子必減少了。
那時，若有人對你們說：基督在這裡，或說：基督在那裡，你們不要信！
因為假基督、假先知將要起來，顯大神蹟、大奇事，倘若能行，連選民也就迷惑了。
看哪，我預先告訴你們了。
若有人對你們說：看哪，基督在曠野裡，你們不要出去！或說：看哪，基督在內屋中，你們不要信！
閃電從東邊發出，直照到西邊。人子降臨，也要這樣。
屍首在哪裡，鷹也必聚在那裡。

\paragraph{馬可 13：14-23}
你們看見那行毀壞可憎的，站在不當站的地方（讀這經的人須要會意）。那時，在猶太的，應當逃到山上；
在房上的，不要下來，也不要進去拿家裡的東西；
在田裡的，也不要回去取衣裳。
當那些日子，懷孕的和奶孩子的有禍了！
你們應當祈求，叫這些事不在冬天臨到。
因為在那些日子必有災難，自從神創造萬物直到如今，並沒有這樣的災難，後來也必沒有。
若不是主減少那日子，凡有血氣的，總沒有一個得救的；只是為主的選民，他將那日子減少了。
那時若有人對你們說：看哪，基督在這裡，或說：基督在那裡，你們不要信！
因為假基督、假先知將要起來，顯神蹟奇事，倘若能行，就把選民迷惑了。
你們要謹慎。看哪，凡事我都預先告訴你們了。

\paragraph{路加 21：20-24}
你們看見耶路撒冷被兵圍困，就可知道他成荒場的日子近了。
那時，在猶太的應當逃到山上；在城裡的應當出來；在鄉下的不要進城；
因為這是報應的日子，使經上所寫的都得應驗。
當那些日子，懷孕的和奶孩子的有禍了！因為將有大災難降在這地方，也有震怒臨到這百姓。
他們要倒在刀下，又被擄到各國去。耶路撒冷要被外邦人踐踏，直到外邦人的日期滿了。
\paragraph{综合}
“当你们看见耶路撒冷被军队包围的时候，你们那时就应当知道它的毁灭近了。「當你們看見藉著先知但以理所說的『那帶來毀滅的褻瀆者 』 站在聖地(不該站的地方)的時候——讀者應當思考——  21 那时，在犹太的人应当逃到山里；在城 里的人应当出来；在乡下的人不要进城；在屋頂上的人，不要下來拿屋子裡的東西； 18 在田裡的人，不要返回去拿自己的衣服。 22 因为这是报应的日子，为要应验经上所记载的一切。 23 在那些日子里，孕妇和哺乳的女人有祸了！你們要禱告，好讓你們逃難的時候不在冬天，也不在安息日。因为将有大苦难临到这地，也有震怒临到这民 h。這樣的患難從世界開始直到如今，從來沒有發生過，也絕不再發生。22 如果那些日子不被主減少，就沒有一個人 f能得救。可是為了那些蒙揀選的人，就是他所揀選的，那些日子將被減少。 24 他们将倒在刀口下，被掳到各国去。耶路撒冷将要被外邦人践踏，直到外邦人的日期满了。
「那時候，如果有人對你們說『看哪，基督在這裡！』或『他在那裡！』你們不要相信。 24 因為假基督們和假先知們會出現，行大神蹟和奇事，如果有可能，甚至迷惑那些蒙揀選的人。 25 看哪，我已經預先告訴你們了！ 26 所以，如果有人對你們說『看哪，基督在曠野裡！』你們不要出去；有人說：『看哪，他在內室裡！』你們也不要相信；27 因為就像閃電從東邊發出，直照到西邊，人子的來臨也將要這樣。 28 屍體在哪裡，禿鷹就聚集在哪裡。


\subsubsection{駕雲降臨-太24:29-35/可13：24-31/路21：25-33}
\paragraph{馬太 24:29-35}
那些日子的災難一過去，日頭就變黑了，月亮也不放光，眾星要從天上墜落，天勢都要震動。
那時，人子的兆頭要顯在天上，地上的萬族都要哀哭。他們要看見人子，有能力，有大榮耀，駕著天上的雲降臨。
他要差遣使者，用號筒的大聲，將他的選民，從四方（方：原文是風），從天這邊到天那邊，都招聚了來。
你們可以從無花果樹學個比方：當樹枝發嫩長葉的時候，你們就知道夏天近了。
這樣，你們看見這一切的事，也該知道人子近了，正在門口了。
我實在告訴你們，這世代還沒有過去，這些事都要成就。
天地要廢去，我的話卻不能廢去。


\paragraph{馬可 13：24-31}
在那些日子，那災難以後，日頭要變黑了，月亮也不放光，
眾星要從天上墜落，天勢都要震動。
那時，他們（馬太二十四章三十節是地上的萬族）要看見人子有大能力、大榮耀，駕雲降臨。
他要差遣天使，把他的選民，從四方（方：原文是風），從地極直到天邊，都招聚了來。
你們可以從無花果樹學個比方：當樹枝發嫩長葉的時候，你們就知道夏天近了。
這樣，你們幾時看見這些事成就，也該知道人子（人子：或作神的國）近了，正在門口了。
我實在告訴你們，這世代還沒有過去，這些事都要成就。
天地要廢去，我的話卻不能廢去。
\paragraph{路加 21：25-33}
日、月、星辰要顯出異兆，地上的邦國也有困苦；因海中波浪的響聲，就慌慌不定。
天勢都要震動，人想起那將要臨到世界的事，就都嚇得魂不附體。
那時，他們要看見人子有能力，有大榮耀駕雲降臨。
一有這些事，你們就當挺身昂首，因為你們得贖的日子近了。
耶穌又設比喻對他們說：你們看無花果樹和各樣的樹；他發芽的時候，你們一看見，自然曉得夏天近了。
這樣，你們看見這些事漸漸的成就，也該曉得神的國近了。
我實在告訴你們，這世代還沒有過去，這些事都要成就。
天地要廢去，我的話卻不能廢去。
\paragraph{综合}
在那些日子裡，隨著那患難，在太阳、月亮和星辰中将有征兆出现；那些日子的患難一過去，『太陽就要變黑，月亮也不發光，星辰要從天上墜落，諸天的各勢力也將被震動。』
在地上，因海洋的咆哮和翻腾，列国的人就惊慌失措，陷入困惑混乱之中； 26 人们预料到将要临到世界的事，就吓得魂不附体；因为诸天的各势力都将被震动。 27 那时候，人子的徵兆將要顯現在天上，地上的萬族都要捶胸哀哭。人们将要看见人子带着极大的权能和荣耀，在云彩中来临。在嘹亮的號角聲中，他要差派他的天使們從四方，從天這邊到天那邊，召集他所揀選的人。 28 这些事一开始发生，你们就当挺起身、抬起头，因为你们的救赎近了。”

\subsubsection{警醒的勸誡-太24:36-51/可13：32-37/路21：34-36}
\paragraph{馬太 24:36-51}
但那日子，那時辰，沒有人知道，連天上的使者也不知道，子也不知道，惟獨父知道。
挪亞的日子怎樣，人子降臨也要怎樣。
當洪水以前的日子，人照常吃喝嫁娶，直到挪亞進方舟的那日；
不知不覺洪水來了，把他們全都沖去。人子降臨也要這樣。
那時，兩個人在田裡，取去一個，撇下一個。
兩個女人推磨，取去一個，撇下一個。
所以，你們要警醒，因為不知道你們的主是哪一天來到。
家主若知道幾更天有賊來，就必警醒，不容人挖透房屋；這是你們所知道的。
所以，你們也要預備，因為你們想不到的時候，人子就來了。
誰是忠心有見識的僕人，為主人所派，管理家裡的人，按時分糧給他們呢？
主人來到，看見他這樣行，那僕人就有福了。
我實在告訴你們，主人要派他管理一切所有的。
倘若那惡僕心裡說：我的主人必來得遲，
就動手打他的同伴，又和酒醉的人一同吃喝。
在想不到的日子，不知道的時辰，那僕人的主人要來，
重重地處治他（或作：把他腰斬了），定他和假冒為善的人同罪；在那裡必要哀哭切齒了。

\paragraph{馬可 13：32-37}
但那日子，那時辰，沒有人知道，連天上的使者也不知道，子也不知道，惟有父知道。
你們要謹慎，警醒祈禱，因為你們不曉得那日期幾時來到。
這事正如一個人離開本家，寄居外邦，把權柄交給僕人，分派各人當做的工，又吩咐看門的警醒。
所以，你們要警醒；因為你們不知道家主什麼時候來，或晚上，或半夜，或雞叫，或早晨；
恐怕他忽然來到，看見你們睡著了。
我對你們所說的話，也是對眾人說：要警醒！

\paragraph{路加 21：34-36}
你們要謹慎，恐怕因貪食、醉酒，並今生的思慮累住你們的心，那日子就如同網羅忽然臨到你們；
因為那日子要這樣臨到全地上一切居住的人。
你們要時時警醒，常常祈求，使你們能逃避這一切要來的事，得以站立在人子面前。
\paragraph{综合}
「你們可以從無花果樹學個比方。當樹枝發嫩長葉的時候，你們就知道夏天近了。29 這樣，你們幾時看見這些事成就，也該知道人子 d近了，正在門口了。 30 我實在告訴你們：這世代還沒有過去，這些事都要成就。 31 天地要廢去，我的話卻不能廢去。32 但那日子、那時辰，沒有人知道，連天上的使者也不知道，子也不知道，唯有父知道。挪亞的那些日子怎樣，人子的來臨也將要怎樣。 38 在洪水以前的那些日子裡，人們繼續吃、喝、嫁、娶，直到挪亞進方舟的那一天。 39 他們毫不察覺，直到洪水到來，把一切沖走。人子的來臨，也將是這樣。 40 那時，兩個人在田裡，一個被接去，一個被留下。 41 兩個女人在磨石那裡推磨，一個被接去，一個被留下。33 你們要謹慎，警醒祈禱，因為你們不曉得那日期幾時來到。你們不知道你們的主哪一天 i回來。 34 這事正如一個人離開本家，寄居外邦，把權柄交給僕人，分派各人當做的工，又吩咐看門的警醒。 35 所以，你們要警醒，因為你們不知道家主什麼時候來，或晚上，或半夜，或雞叫，或早晨。 36 恐怕他忽然來到，看見你們睡著了。不過你們應當知道這一點：一家的主人如果知道賊什麼時刻 j來，他就會警醒，不讓他的房子被人鑽進。 44 你們為了這緣故也應當做好準備，因為在你們意想不到的時候，人子就來了。 37 我對你們所說的話，也是對眾人說：要警醒！」“你们要谨慎自守，免得你们的心因宴乐 l、醉酒和今生的忧虑变得迟钝，那日子就会像网罗一样突然临到你们， 35 因为那日子将要临到所有住在全地面上的人。 36 你们要时刻警醒、祷告，使你们能够逃避这一切将要发生的事，好站立在人子面前。”
「究竟誰是那又忠心又聰明的奴僕——受主人委任統管他家裡的人、按時給他們分糧的人呢？ 46 主人回來的時候，看見哪個奴僕這樣做，那個奴僕就蒙福了。 47 我確實地告訴你們：主人會委任他統管自己所擁有的一切。 48 但如果這惡奴心裡說『我的主人會遲延回來 k』， 49 就動手毆打與他同做奴僕的，並且與醉酒的人一起吃喝， 50 那麼，這奴僕的主人就要在意想不到的日子、在他不知道的時刻回來， 51 並且會嚴厲懲罰他 l，使他與那些偽善的人有同樣的下場。在那裡將有哀哭和切齒。


\subsubsection{十個童女-馬太 25:1-13}
那時，天國好比十個童女拿著燈出去迎接新郎。
其中有五個是愚拙的，五個是聰明的。
愚拙的拿著燈，卻不預備油；
聰明的拿著燈，又預備油在器皿裡。
新郎遲延的時候，他們都打盹，睡著了。
半夜有人喊著說：新郎來了，你們出來迎接他！
那些童女就都起來收拾燈。
愚拙的對聰明的說：請分點油給我們，因為我們的燈要滅了。
聰明的回答說：恐怕不夠你我用的；不如你們自己到賣油的那裡去買吧。
他們去買的時候，新郎到了。那預備好了的，同他進去坐席，門就關了。
其餘的童女隨後也來了，說：主啊，主啊，給我們開門！
他卻回答說：我實在告訴你們，我不認識你們。
所以，你們要警醒；因為那日子，那時辰，你們不知道。

\subsubsection{按才幹給銀子-馬太 25:14-30}
天國又好比一個人要往外國去，就叫了僕人來，把他的家業交給他們，
按著各人的才幹給他們銀子：一個給了五千，一個給了二千，一個給了一千，就往外國去了。
那領五千的隨即拿去做買賣，另外賺了五千。
那領二千的也照樣另賺了二千。
但那領一千的去掘開地，把主人的銀子埋藏了。
過了許久，那些僕人的主人來了，和他們算帳。
那領五千銀子的又帶著那另外的五千來，說：主啊，你交給我五千銀子。請看，我又賺了五千。
主人說：好，你這又良善又忠心的僕人，你在不多的事上有忠心，我要把許多事派你管理；可以進來享受你主人的快樂。
那領二千的也來，說：主啊，你交給我二千銀子。請看，我又賺了二千。
主人說：好，你這又良善又忠心的僕人，你在不多的事上有忠心，我要把許多事派你管理；可以進來享受你主人的快樂。
那領一千的也來，說：主啊，我知道你是忍心的人，沒有種的地方要收割，沒有散的地方要聚斂，
我就害怕，去把你的一千銀子埋藏在地裡。請看，你的原銀子在這裡。
主人回答說：你這又惡又懶的僕人，你既知道我沒有種的地方要收割，沒有散的地方要聚斂，
就當把我的銀子放給兌換銀錢的人，到我來的時候，可以連本帶利收回。
奪過他這一千來，給那有一萬的。
因為凡有的，還要加給他，叫他有餘；沒有的，連他所有的也要奪過來。
把這無用的僕人丟在外面黑暗裡；在那裡必要哀哭切齒了。

\subsubsection{分別綿羊山羊-馬太25:31-46}
當人子在他榮耀裡、同著眾天使降臨的時候，要坐在他榮耀的寶座上。
萬民都要聚集在他面前。他要把他們分別出來，好像牧羊的分別綿羊山羊一般，
把綿羊安置在右邊，山羊在左邊。
於是王要向那右邊的說：你們這蒙我父賜福的，可來承受那創世以來為你們所預備的國；
因為我餓了，你們給我吃，渴了，你們給我喝；我作客旅，你們留我住；
我赤身露體，你們給我穿；我病了、你們看顧我；我在監裡，你們來看我。
義人就回答說：主啊，我們什麼時候見你餓了，給你吃，渴了，給你喝？
什麼時候見你作客旅，留你住，或是赤身露體，給你穿？
又什麼時候見你病了，或是在監裡，來看你呢？
王要回答說：我實在告訴你們，這些事你們既做在我這弟兄中一個最小的身上，就是做在我身上了。
王又要向那左邊的說：你們這被咒詛的人，離開我！進入那為魔鬼和他的使者所預備的永火裡去！
因為我餓了，你們不給我吃，渴了，你們不給我喝；
我作客旅，你們不留我住；我赤身露體，你們不給我穿；我病了，我在監裡，你們不來看顧我。
他們也要回答說：主啊，我們什麼時候見你餓了，或渴了，或作客旅，或赤身露體，或病了，或在監裡，不伺候你呢？
王要回答說：我實在告訴你們，這些事你們既不做在我這弟兄中一個最小的身上，就是不做在我身上了。
這些人要往永刑裡去；那些義人要往永生裡去。



%\subsubsection{約翰}

\subsection{門徒預備逾越節筵席}

\subsubsection{馬太 26：17-19}
除酵節的第一天，門徒來問耶穌說：你吃逾越節的筵席，要我們在哪裡給你預備？
耶穌說：你們進城去，到某人那裡，對他說：夫子說：我的時候快到了，我與門徒要在你家裡守逾越節。
門徒遵著耶穌所吩咐的就去預備了逾越節的筵席。
\subsubsection{馬可 14：12-16}
除酵節的第一天，就是宰逾越羊羔的那一天，門徒對耶穌說：你吃逾越節的筵席要我們往哪裡去預備呢？
耶穌就打發兩個門徒，對他們說：你們進城去，必有人拿著一瓶水迎面而來，你們就跟著他。
他進哪家去，你們就對那家的主人說：夫子說：客房在哪裡？我與門徒好在那裡吃逾越節的筵席。
他必指給你們擺設整齊的一間大樓，你們就在那裡為我們預備。
門徒出去，進了城，所遇見的正如耶穌所說的。他們就預備了逾越節的筵席。
\subsubsection{路加 22：7-13}
除酵節，須宰逾越羊羔的那一天到了。
耶穌打發彼得、約翰，說：你們去為我們預備逾越節的筵席，好叫我們吃。
他們問他說：要我們在哪裡預備？
耶穌說：你們進了城，必有人拿著一瓶水迎面而來，你們就跟著他，到他所進的房子裡去，
對那家的主人說：夫子說：客房在哪裡？我與門徒好在那裡吃逾越節的筵席。
他必指給你們擺設整齊的一間大樓，你們就在那裡預備。
他們去了，所遇見的正如耶穌所說的；他們就預備了逾越節的筵席。
\subsubsection{综合}
除酵節的第一天，就是宰逾越羊羔的那一天，門徒對耶穌說：「你吃逾越節的筵席，要我們往哪裡去預備呢？」耶穌就打發兩個門徒彼得和约翰，對他們說：你們去為我們預備逾越節的筵席，好叫我們吃。」他們問他說：「要我們在哪裡預備？」 耶穌說：「你們進城去，必有人拿著一瓶水迎面而來，你們就跟著他。他進哪家去，你們就對那家的主人說：『夫子說：客房在哪裡？我的時候快到了，我與門徒好在那裡吃逾越節的筵席。』他必指給你們擺設整齊的一間大樓，你們就在那裡為我們預備。」門徒出去，進了城，所遇見的正如耶穌所說的，他們就預備了逾越節的筵席。

\subsection{逾越節筵席}


\subsubsection{設立聖餐-太26：26-29/可14：22-26/路22：14-20}
\paragraph{馬太 26：26-29}
他們吃的時候，耶穌拿起餅來，祝福，就擘開，遞給門徒，說：你們拿著吃，這是我的身體；
又拿起杯來，祝謝了，遞給他們，說：你們都喝這個；
因為這是我立約的血，為多人流出來，使罪得赦。
但我告訴你們，從今以後，我不再喝這葡萄汁，直到我在我父的國裡同你們喝新的那日子。
\paragraph{馬可 14：22-26}
他們吃的時候，耶穌拿起餅來，祝了福，就擘開，遞給他們，說：你們拿著吃，這是我的身體；
又拿起杯來，祝謝了，遞給他們；他們都喝了。
耶穌說：這是我立約的血，為多人流出來的。
我實在告訴你們，我不再喝這葡萄汁，直到我在神的國裡喝新的那日子。
\paragraph{路加 22：14-20}
時候到了，耶穌坐席，使徒也和他同坐。
耶穌對他們說：我很願意在受害以先和你們吃這逾越節的筵席。
我告訴你們，我不再吃這筵席，直到成就在神的國裡。
耶穌接過杯來，祝謝了，說：你們拿這個，大家分著喝。
我告訴你們，從今以後，我不再喝這葡萄汁，直等神的國來到。
又拿起餅來，祝謝了，就擘開，遞給他們，說：這是我的身體，為你們捨的，你們也應當如此行，為的是記念我。
飯後也照樣拿起杯來，說：這杯是用我血所立的新約，是為你們流出來的。
\paragraph{综合}
到了晚上，耶穌和十二個門徒都來了。時候到了，耶穌坐席，使徒也和他同坐。耶穌對他們說：「我很願意在受害以先和你們吃這逾越節的筵席。我告訴你們，我不再吃這筵席，直到成就在神的國裡。」他們吃的時候，耶穌拿起餅來，祝福，就掰開，遞給門徒，說：「你們拿著吃，這是我的身體。為你們捨的。你們也應當如此行，為的是記念我。」飯後又拿起杯來，祝謝了，遞給他們，說：「你們都喝這個，大家分著喝。因為這是我立約的血，為多人流出來，使罪得赦。但我實在告訴你們：從今以後，我不再喝這葡萄汁，直到我在我父的國裡同你們喝新的那日子。」他們都喝了。




\subsubsection{耶穌預言猶大賣主-太26：20-25/可14：17-21/路22：21-22/約13：18-30}
\paragraph{馬太 26：20-25}
到了晚上，耶穌和十二個門徒坐席。
正吃的時候，耶穌說：我實在告訴你們，你們中間有一個人要賣我了。
他們就甚憂愁，一個一個的問他說：主，是我嗎？
耶穌回答說：同我蘸手在盤子裡的，就是他要賣我。
人子必要去世，正如經上指著他所寫的；但賣人子的人有禍了！那人不生在世上倒好。
賣耶穌的猶大問他說：拉比，是我嗎？耶穌說：你說的是。
\paragraph{馬可 14：17-21}
到了晚上，耶穌和十二個門徒都來了。
他們坐席正吃的時候，耶穌說：我實在告訴你們，你們中間有一個與我同吃的人要賣我了。
他們就憂愁起來，一個一個的問他說：是我嗎？
耶穌對他們說：是十二個門徒中同我蘸手在盤子裡的那個人。
人子必要去世，正如經上指著他所寫的；但賣人子的人有禍了！那人不生在世上倒好。
\paragraph{路加 22：21-23}
看哪！那賣我之人的手與我一同在桌子上。
人子固然要照所預定的去世，但賣人子的人有禍了！
他們就彼此對問，是哪一個要做這事。
\paragraph{約翰 13：18-30}
我這話不是指著你們眾人說的，我知道我所揀選的是誰。現在要應驗經上的話，說：同我吃飯的人用腳踢我。
如今事情還沒有成就，我要先告訴你們，叫你們到事情成就的時候可以信我是基督。
我實實在在的告訴你們，有人接待我所差遣的，就是接待我；接待我，就是接待那差遣我的。
耶穌說了這話，心裡憂愁，就明說：我實實在在的告訴你們，你們中間有一個人要賣我了。
門徒彼此對看，猜不透所說的是誰。
有一個門徒，是耶穌所愛的，側身挨近耶穌的懷裡。
西門彼得點頭對他說：你告訴我們，主是指著誰說的。
那門徒便就勢靠著耶穌的胸膛，問他說：主啊，是誰呢？
耶穌回答說：我蘸一點餅給誰，就是誰。耶穌就蘸了一點餅，遞給加略人西門的兒子猶大。
他吃了以後，撒但就入了他的心。耶穌便對他說：你所做的，快做吧！
同席的人沒有一個知道是為什麼對他說這話。
有人因猶大帶著錢囊，以為耶穌是對他說：你去買我們過節所應用的東西，或是叫他拿什麼賙濟窮人。
猶大受了那點餅，立刻就出去。那時候是夜間了。
\paragraph{综合}
正吃的時候，耶穌心裡憂愁說：「我實在告訴你們，你們中間有一個我同吃的人要賣我了。」門徒彼此對看，猜不透所說的是誰。他們就甚憂愁，一個一個的問他說：「主，是我嗎？」耶穌回答說：「是十二個門徒中同我蘸手在盤子裡的，就是他要賣我。 人子必要去世，正如經上指著他所寫的；但賣人子的人有禍了！那人不生在世上倒好。」。 23 有一個門徒，是耶穌所愛的，側身挨近耶穌的懷裡。 24 西門彼得點頭對他說：「你告訴我們主是指著誰說的。」 25 那門徒便就勢靠著耶穌的胸膛，問他說：「主啊，是誰呢？」 26 耶穌回答說：「我蘸一點餅給誰，就是誰。」耶穌就蘸了一點餅，遞給加略人西門的兒子猶大。賣耶穌的猶大問他說：「拉比，是我嗎？」耶穌說：「你說的是。」他吃了以後，撒旦就入了他的心。耶穌便對他說：「你所做的，快做吧！」 28 同席的人沒有一個知道為什麼對他說這話。 29 有人因猶大帶著錢囊，以為耶穌是對他說「你去買我們過節所應用的東西」，或是叫他拿什麼賙濟窮人。 30 猶大受了那點餅，立刻就出去。那時候是夜間了。 

\subsubsection{筵席上門徒爭大小-路加 22：24-30}

門徒起了爭論，他們中間哪一個可算為大。
耶穌說：外邦人有君王為主治理他們，那掌權管他們的稱為恩主。
但你們不可這樣；你們裡頭為大的，倒要像年幼的；為首領的，倒要像服事人的。
是誰為大？是坐席的呢？是服事人的呢？不是坐席的大嗎？然而，我在你們中間如同服事人的。
我在磨煉之中，常和我同在的就是你們。
我將國賜給你們，正如我父賜給我一樣，
叫你們在我國裡，坐在我的席上吃喝，並且坐在寶座上，審判以色列十二個支派。


\subsubsection{耶穌給門徒洗脚-約翰13：1-17}

逾越節以前，耶穌知道自己離世歸父的時候到了。他既然愛世間屬自己的人，就愛他們到底。
吃晚飯的時候，魔鬼已將賣耶穌的意思放在西門的兒子加略人猶大心裡。
耶穌知道父已將萬有交在他手裡，且知道自己是從神出來的，又要歸到神那裡去，
就離席站起來，脫了衣服，拿一條手巾束腰，
隨後把水倒在盆裡，就洗門徒的腳，並用自己所束的手巾擦乾。
挨到西門彼得，彼得對他說：主啊，你洗我的腳嗎？
耶穌回答說：我所做的，你如今不知道，後來必明白。
彼得說：你永不可洗我的腳！耶穌說：我若不洗你，你就與我無分了。
西門彼得說：主啊，不但我的腳，連手和頭也要洗。
耶穌說：凡洗過澡的人，只要把腳一洗，全身就乾淨了。你們是乾淨的，然而不都是乾淨的。
耶穌原知道要賣他的是誰，所以說：你們不都是乾淨的。
耶穌洗完了他們的腳，就穿上衣服，又坐下，對他們說：我向你們所做的，你們明白嗎？
你們稱呼我夫子，稱呼我主，你們說的不錯，我本來是。
我是你們的主，你們的夫子，尚且洗你們的腳，你們也當彼此洗腳。
我給你們作了榜樣，叫你們照著我向你們所做的去做。
我實實在在的告訴你們，僕人不能大於主人，差人也不能大於差他的人。
你們既知道這事，若是去行就有福了。
\subsubsection{彼此相愛的命令-約翰13：31-35}

他既出去，耶穌就說：如今人子得了榮耀，神在人子身上也得了榮耀。
神要因自己榮耀人子，並且要快快的榮耀他。
小子們，我還有不多的時候與你們同在；後來你們要找我，但我所去的地方你們不能到。這話我曾對猶太人說過，如今也照樣對你們說。
我賜給你們一條新命令，乃是叫你們彼此相愛；我怎樣愛你們，你們也要怎樣相愛。
你們若有彼此相愛的心，眾人因此就認出你們是我的門徒了。
\subsubsection{	首次預言彼得不認主-可14:27-31/路22:31-34/約13:36-38}
\paragraph{馬可 14：27-31}
耶穌對他們說：你們都要跌倒了，因為經上記著說：我要擊打牧人，羊就分散了。
但我復活以後，要在你們以先往加利利去。
彼得說：眾人雖然跌倒，我總不能。
耶穌對他說：我實在告訴你，就在今天夜裡，雞叫兩遍以先，你要三次不認我。
彼得卻極力的說：我就是必須和你同死，也總不能不認你。眾門徒都是這樣說。
\paragraph{路加 22：31-34}
主又說：西門！西門！撒但想要得著你們，好篩你們像篩麥子一樣；
但我已經為你祈求，叫你不至於失了信心，你回頭以後，要堅固你的弟兄。
彼得說：主啊，我就是同你下監，同你受死，也是甘心！
耶穌說：彼得，我告訴你，今日雞還沒有叫，你要三次說不認得我。
\paragraph{約翰 13：36-38}
西門彼得問耶穌說：主往哪裡去？耶穌回答說：我所去的地方，你現在不能跟我去，後來卻要跟我去。
彼得說：主啊，我為什麼現在不能跟你去？我願意為你捨命！
耶穌說：你願意為我捨命嗎？我實實在在的告訴你，雞叫以先，你要三次不認我。
\paragraph{综合}
西門彼得問耶穌說：「主往哪裡去？」耶穌回答說：「我所去的地方，你現在不能跟我去，後來卻要跟我去。」彼得說：「主啊，我為什麼現在不能跟你去？我願意為你捨命！」耶穌說：「你願意為我捨命嗎？我實實在在地告訴你：雞叫以先，你要三次不認我。」主又說：「西門，西門，撒旦想要得著你們，好篩你們像篩麥子一樣。但我已經為你祈求，叫你不至於失了信心。你回頭以後，要堅固你的弟兄。」彼得說：「主啊，我就是同你下監，同你受死，也是甘心！」耶穌說：「彼得，我告訴你：今日雞還沒有叫，你要三次說不認得我。」耶穌又對他們說：「我差你們出去的時候，沒有錢囊、沒有口袋、沒有鞋，你們缺少什麼沒有？」他們說：「沒有。」耶穌說：「但如今有錢囊的可以帶著，有口袋的也可以帶著，沒有刀的要賣衣服買刀。我告訴你們，經上寫著說『他被列在罪犯之中』，這話必應驗在我身上，因為那關係我的事必然成就。」他們說：「主啊，請看，這裡有兩把刀！」耶穌說：「夠了。」


\subsubsection{耶穌囑咐門徒外出的原則-路22：35-38}

%\paragraph{路加 22：36-38}
耶穌又對他們說：我差你們出去的時候，沒有錢囊，沒有口袋，沒有鞋，你們缺少什麼沒有？他們說：沒有。
耶穌說：但如今有錢囊的可以帶著，有口袋的也可以帶著，沒有刀的要賣衣服買刀。
我告訴你們，經上寫著說：他被列在罪犯之中。這話必應驗在我身上；因為那關係我的事必然成就。
他們說：主啊，請看！這裡有兩把刀。耶穌說：夠了。

\subsubsection{耶穌告訴門徒自己去預備地方-約翰14：1-4}
%\paragraph{約翰}
你們心裡不要憂愁；你們信神，也當信我。
在我父的家裡有許多住處；若是沒有，我就早已告訴你們了。我去原是為你們預備地方去。
我若去為你們預備了地方，就必再來接你們到我那裡去，我在哪裡，叫你們也在那裡。
我往哪裡去，你們知道；那條路，你們也知道（有古卷作：我往哪裡去，你們知道那條路）。
\subsubsection{回答多馬的問題-約翰14：5-7}
多馬對他說：主啊，我們不知道你往哪裡去，怎麼知道那條路呢？
耶穌說我就是道路、真理、生命；若不藉著我，沒有人能到父那裡去。
你們若認識我，也就認識我的父。從今以後，你們認識他，並且已經看見他。
\subsubsection{回答腓力的問題-約翰14：8-21}
腓力對他說：求主將父顯給我們看，我們就知足了。
耶穌對他說：腓力，我與你們同在這樣長久，你還不認識我嗎？人看見了我，就是看見了父；你怎麼說將父顯給我們看呢？
我在父裡面，父在我裡面，你不信嗎？我對你們所說的話，不是憑著自己說的，乃是住在我裡面的父做他自己的事。
你們當信我，我在父裡面，父在我裡面；即或不信，也當因我所做的事信我。
我實實在在的告訴你們，我所做的事，信我的人也要做，並且要做比這更大的事，因為我往父那裡去。
你們奉我的名無論求什麼，我必成就，叫父因兒子得榮耀。
你們若奉我的名求什麼，我必成就。
你們若愛我，就必遵守我的命令。
我要求父，父就另外賜給你們一位保惠師（或作：訓慰師；下同），叫他永遠與你們同在，
就是真理的聖靈，乃世人不能接受的；因為不見他，也不認識他。你們卻認識他，因他常與你們同在，也要在你們裡面。
我不撇下你們為孤兒，我必到你們這裡來。
還有不多的時候，世人不再看見我，你們卻看見我；因為我活著，你們也要活著。
到那日，你們就知道我在父裡面，你們在我裡面，我也在你們裡面。
有了我的命令又遵守的，這人就是愛我的；愛我的必蒙我父愛他，我也要愛他，並且要向他顯現。
\subsubsection{回答猶大的問題-約翰14：22-31}
猶大（不是加略人猶大）問耶穌說：主啊，為什麼要向我們顯現，不向世人顯現呢？
耶穌回答說：人若愛我，就必遵守我的道；我父也必愛他，並且我們要到他那裡去，與他同住。
不愛我的人就不遵守我的道。你們所聽見的道不是我的，乃是差我來之父的道。
我還與你們同住的時候，已將這些話對你們說了。
但保惠師，就是父因我的名所要差來的聖靈，他要將一切的事指教你們，並且要叫你們想起我對你們所說的一切話。
我留下平安給你們；我將我的平安賜給你們。我所賜的，不像世人所賜的。你們心裡不要憂愁，也不要膽怯。
你們聽見我對你們說了，我去還要到你們這裡來。你們若愛我，因我到父那裡去，就必喜樂，因為父是比我大的。
現在事情還沒有成就，我預先告訴你們，叫你們到事情成就的時候就可以信。
以後我不再和你們多說話，因為這世界的王將到。他在我裡面是毫無所有；
但要叫世人知道我愛父，並且父怎樣吩咐我，我就怎樣行。起來，我們走吧！

\subsubsection{有關聖靈的教導-約翰15：1-16：16}
我是真葡萄樹，我父是栽培的人。
凡屬我不結果子的枝子，他就剪去；凡結果子的，他就修理乾淨，使枝子結果子更多。
現在你們因我講給你們的道，已經乾淨了。
你們要常在我裡面，我也常在你們裡面。枝子若不常在葡萄樹上，自己就不能結果子；你們若不常在我裡面，也是這樣。
我是葡萄樹，你們是枝子。常在我裡面的，我也常在他裡面，這人就多結果子；因為離了我，你們就不能做什麼。
人若不常在我裡面，就像枝子丟在外面枯乾，人拾起來，扔在火裡燒了。
你們若常在我裡面，我的話也常在你們裡面，凡你們所願意的，祈求，就給你們成就。
你們多結果子，我父就因此得榮耀，你們也就是我的門徒了。
我愛你們，正如父愛我一樣；你們要常在我的愛裡。
你們若遵守我的命令，就常在我的愛裡，正如我遵守了我父的命令，常在他的愛裡。
這些事我已經對你們說了，是要叫我的喜樂存在你們心裡，並叫你們的喜樂可以滿足。
你們要彼此相愛，像我愛你們一樣；這就是我的命令。
人為朋友捨命，人的愛心沒有比這個大的。
你們若遵行我所吩咐的，就是我的朋友了。
以後我不再稱你們為僕人，因僕人不知道主人所做的事。我乃稱你們為朋友；因我從我父所聽見的，已經都告訴你們了。
不是你們揀選了我，是我揀選了你們，並且分派你們去結果子，叫你們的果子常存，使你們奉我的名，無論向父求什麼，他就賜給你們。
我這樣吩咐你們，是要叫你們彼此相愛。
世人若恨你們，你們知道（或作：該知道），恨你們以先已經恨我了。
你們若屬世界，世界必愛屬自己的；只因你們不屬世界，乃是我從世界中揀選了你們，所以世界就恨你們。
你們要記念我從前對你們所說的話：僕人不能大於主人。他們若逼迫了我，也要逼迫你們；若遵守了我的話，也要遵守你們的話。
但他們因我的名要向你們行這一切的事，因為他們不認識那差我來的。
我若沒有來教訓他們，他們就沒有罪；但如今他們的罪無可推諉了。
恨我的，也恨我的父。
我若沒有在他們中間行過別人未曾行的事，他們就沒有罪；但如今連我與我的父，他們也看見也恨惡了。
這要應驗他們律法上所寫的話，說：他們無故的恨我。
但我要從父那裡差保惠師來，就是從父出來真理的聖靈；他來了，就要為我作見證。
你們也要作見證，因為你們從起頭就與我同在。



我已將這些事告訴你們，使你們不至於跌倒。
人要把你們趕出會堂，並且時候將到，凡殺你們的就以為是事奉神。
他們這樣行，是因未曾認識父，也未曾認識我。
我將這事告訴你們，是叫你們到了時候可以想起我對你們說過了。我起先沒有將這事告訴你們，因為我與你們同在。
現今我往差我來的父那裡去，你們中間並沒有人問我：你往哪裡去？
只因我將這事告訴你們，你們就滿心憂愁。
然而，我將真情告訴你們，我去是於你們有益的；我若不去，保惠師就不到你們這裡來；我若去，就差他來。
他既來了，就要叫世人為罪、為義、為審判，自己責備自己。
為罪，是因他們不信我；
為義，是因我往父那裡去，你們就不再見我；
為審判，是因這世界的王受了審判。
我還有好些事要告訴你們，但你們現在擔當不了（或作：不能領會）。
只等真理的聖靈來了，他要引導你們明白（原文作進入）一切的真理；因為他不是憑自己說的，乃是把他所聽見的都說出來，並要把將來的事告訴你們。
他要榮耀我，因為他要將受於我的告訴你們。
凡父所有的，都是我的；所以我說，他要將受於我的告訴你們。
等不多時，你們就不得見我；再等不多時，你們還要見我。
\subsubsection{回答幾個門徒的問題-約翰16：17-33}
有幾個門徒就彼此說：他對我們說：等不多時，你們就不得見我；再等不多時，你們還要見我；又說：因我往父那裡去。這是什麼意思呢？
門徒彼此說：他說等不多時到底是什麼意思呢？我們不明白他所說的話。
耶穌看出他們要問他，就說：我說等不多時，你們就不得見我；再等不多時，你們還要見我，你們為這話彼此相問嗎？
我實實在在的告訴你們，你們將要痛哭、哀號，世人倒要喜樂；你們將要憂愁，然而你們的憂愁要變為喜樂。
婦人生產的時候就憂愁，因為他的時候到了；既生了孩子，就不再記念那苦楚，因為歡喜世上生了一個人。
你們現在也是憂愁，但我要再見你們，你們的心就喜樂了；這喜樂也沒有人能奪去。
到那日，你們什麼也就不問我了。我實實在在的告訴你們，你們若向父求什麼，他必因我的名賜給你們。
向來你們沒有奉我的名求什麼，如今你們求，就必得著，叫你們的喜樂可以滿足。
這些事，我是用比喻對你們說的；時候將到，我不再用比喻對你們說，乃要將父明明的告訴你們。
到那日，你們要奉我的名祈求；我並不對你們說，我要為你們求父。
父自己愛你們；因為你們已經愛我，又信我是從父出來的。
我從父出來，到了世界；我又離開世界，往父那裡去。
門徒說：如今你是明說，並不用比喻了。
現在我們曉得你凡事都知道，也不用人問你，因此我們信你是從神出來的。
耶穌說：現在你們信嗎？
看哪，時候將到，且是已經到了，你們要分散，各歸自己的地方去，留下我獨自一人；其實我不是獨自一人，因為有父與我同在。
我將這些事告訴你們，是要叫你們在我裡面有平安。在世上，你們有苦難；但你們可以放心，我已經勝了世界。






\subsubsection{耶穌為自己的禱告-約翰17：1-8}
耶穌說了這話，就舉目望天，說：父啊，時候到了，願你榮耀你的兒子，使兒子也榮耀你；正如你曾賜給他權柄管理凡有血氣的，叫他將永生賜給你所賜給他的人。認識你─獨一的真神，並且認識你所差來的耶穌基督，這就是永生。
我在地上已經榮耀你，你所託付我的事，我已成全了。
父啊，現在求你使我同你享榮耀，就是未有世界以先，我同你所有的榮耀。
你從世上賜給我的人，我已將你的名顯明與他們。他們本是你的，你將他們賜給我，他們也遵守了你的道。
如今他們知道，凡你所賜給我的，都是從你那裡來的；
因為你所賜給我的道，我已經賜給他們，他們也領受了，又確實知道，我是從你出來的，並且信你差了我來。

\subsubsection{耶穌為門徒的禱告-約翰17：9-19a}
我為他們祈求，不為世人祈求，卻為你所賜給我的人祈求，因他們本是你的。
凡是我的，都是你的；你的也是我的，並且我因他們得了榮耀。
從今以後，我不在世上，他們卻在世上；我往你那裡去。聖父啊，求你因你所賜給我的名保守他們，叫他們合而為一像我們一樣。
我與他們同在的時候，因你所賜給我的名保守了他們，我也護衛了他們；其中除了那滅亡之子，沒有一個滅亡的，好叫經上的話得應驗。
現在我往你那裡去，我還在世上說這話，是叫他們心裡充滿我的喜樂。
我已將你的道賜給他們。世界又恨他們；因為他們不屬世界，正如我不屬世界一樣。
我不求你叫他們離開世界，只求你保守他們脫離那惡者（或作：脫離罪惡）。
他們不屬世界，正如我不屬世界一樣。
求你用真理使他們成聖；你的道就是真理。
你怎樣差我到世上，我也照樣差他們到世上。
我為他們的緣故，自己分別為聖，叫他們也因真理成聖。

\subsubsection{耶穌為一切信的人禱告-約翰17：20-26}
我不但為這些人祈求，也為那些因他們的話信我的人祈求，
使他們都合而為一。正如你父在我裡面，我在你裡面，使他們也在我們裡面，叫世人可以信你差了我來。
你所賜給我的榮耀，我已賜給他們，使他們合而為一，像我們合而為一。
我在他們裡面，你在我裡面，使他們完完全全的合而為一，叫世人知道你差了我來，也知道你愛他們如同愛我一樣。
父啊，我在哪裡，願你所賜給我的人也同我在那裡，叫他們看見你所賜給我的榮耀；因為創立世界以前，你已經愛我了。
公義的父啊，世人未曾認識你，我卻認識你；這些人也知道你差了我來。
我已將你的名指示他們，還要指示他們，使你所愛我的愛在他們裡面，我也在他們裡面。

















